%% kapa: exam=1 round=27 year=2016 period=2 form=A subject=accounting subject_ko=회계학 seq=2 number=42
%% source: Q-Net 감정평가사 기출문제 (cst00309) artlSeq=5169402 file=extracted/1008/5169402/제27회 감정평가사 1차 문제/제27회 감정평가사 2교시문제(A형).pdf
%% conversion: kordoc 4.24.51 -> tools/gen1.py (auto); review=proofread-sonnet; flags=none
\begin{question}{42}
(주)감평은 20×5년 초 액면금액 ￦1,000,000(액면이자율 연 4\%, 매년 말 이자 지급, 만기 3년)의 전환사채를 발행하였다. 사채 액면금액 ￦3,000당 보통주 (액면금액 ￦1,000) 1주로 전환할 수 있는 권리가 부여되어 있다. 만약 만기일까지 전환권이 행사되지 않을 경우 추가로 ￦198,600의 상환할증금을 지급한다. 이 사채는 액면금액인 ￦1,000,000에 발행되었으며 전환권이 없었다면 ￦949,213에 발행되었을 것이다(유효이자율 연 12\%). 사채발행일 후 1년 된 시점인 20×6년 초에 액면금액의 60\%에 해당하는 전환사채가 보통주로 전환되었다. 이러한 전환으로 인해 증가할 주식발행초과금은? (단, 전환사채 발행 시 인식한 전환권대가 중 전환된 부분은 주식발행초과금으로 대체하며, 단수 차이가 있으면 가장 근사치를 선택한다.)
\begin{choices}
  \item ￦413,871
  \item ￦433,871
  \item ￦444,071
  \item ￦444,343
  \item ￦464,658
\end{choices}
\end{question}
\answer{4}
